\question[4]
Ein Anbieter einer kostenlosen Fitness-App speichert dauerhaft Standort, Puls, Schlafzeiten und
Bewegungsprofile aller Nutzerinnen und Nutzer und wertet diese zentral aus. \\
a. Nenne zwei Gründe, warum die massenhafte Erhebung und Speicherung dieser Daten für die
Betroffenen problematisch sein kann. \\
b. Beurteile das Geschäftsmodell \enquote{kostenlose App gegen Daten} und nenne zwei Maßnahmen,
mit denen die Datenverarbeitung datensparsamer gestaltet werden könnte.

\begin{solutionbox}{6cm}
a. Zum Beispiel: Aus den Bewegungsprofilen lassen sich Wohnort, Arbeitsplatz und Tagesablauf
ableiten (Rückschlüsse auf einzelne Personen trotz \enquote{anonymer} Speicherung). Gesundheitsdaten
sind besonders sensibel und können bei Weitergabe an Dritte (Versicherungen, Arbeitgeber) oder bei
einem Datenleck missbraucht werden. Die Daten können mit anderen Datenquellen zusammengeführt
werden; eine spätere Zweckänderung ist für die Betroffenen kaum kontrollierbar. \\
b. Die Nutzung ist nur scheinbar kostenlos -- bezahlt wird mit persönlichen Daten, deren Wert und
spätere Verwendung die Nutzenden nicht überblicken. Datensparsamer wäre z.\,B.: nur die wirklich
benötigten Daten erheben, Auswertung lokal auf dem Gerät statt zentral, frühzeitige Löschung bzw.
Speicherung nur in aggregierter Form, echte Einwilligung mit klarer Zweckbindung.
\end{solutionbox}
